% Corrected French for printed p.155.
% Corrections: two elisions and the omitted bar in the sheafification target.
\SourcePage{155}

\LayerTarget{tohoku:construction:sheafification:p2}
\noindent \textit{foncteur} en $F$, et l'homomorphisme
$F\to\mathfrak F(\overline F)$ est fonctoriel). Si on veut
interpréter la notion de faisceau de groupes (ou de faisceau de groupes
abéliens etc.) en termes d'espaces étalés, il faut se donner sur chaque
``fibre'' de l'espace étalé une loi de groupe (resp. de groupe abélien
etc.) de façon à satisfaire à une condition naturelle de continuité : on
retrouve alors la définition de [4, XIV].
\LayerTarget{tohoku:remark:stalkwise-exactness}
--- On voit immédiatement sur les espaces étalés correspondants quand un
homomorphisme $F\to G$ de faisceaux est un monomorphisme resp. un
épimorphisme : il faut et il suffit que sur chaque ``fibre'' $F(x)$,
l'homomorphisme correspondant $F(x)\to G(x)$ soit un
monomorphisme resp. un épimorphisme. De même, la fibre en $x$ du noyau,
conoyau, image, coimage d'un homomorphisme de faisceaux de groupes abéliens
s'obtient en prenant le noyau, conoyau etc. de l'homomorphisme de groupes
abéliens $F(x)\to G(x)$.

\LayerTarget{tohoku:def:sheaf-of-modules}
Soit $\mathbf O$ un faisceau d'anneaux avec unité sur $X$, un
\textit{faisceau de $\mathbf O$-modules à gauche} ou bref un
\textit{$\mathbf O$-Module à gauche} sur $X$ est un faisceau de groupes abéliens $F$,
avec la donnée pour tout ouvert $U\subset X$ d'une structure de
$\mathbf O(U)$-module unitaire à gauche sur $F(U)$, de telle façon que les
opérations de restrictions $F(U)\to F(V)$ soient compatibles
avec les opérations de $\mathbf O(U)$ et $\mathbf O(V)$ sur $F(U)$ resp.
$F(V)$. (On peut exprimer ceci aussi en disant qu'on s'est donné un
homomorphisme du faisceau d'anneaux $\mathbf O$ dans le faisceau d'anneaux des
germes d'endomorphismes du faisceau de groupes abéliens $F$).
On définit de façon analogue les $\mathbf O$-Modules à droite ; pour abréger,
nous dirons simplement $\mathbf O$-Modules au lieu de $\mathbf O$-Modules à
gauche. La notion d'homomorphisme de $\mathbf O$-Modules, et de la
composition et addition de tels homomorphismes est évidente, on obtient
alors, \textit{la catégorie additive des $\mathbf O$-Modules sur $X$}, notée
$\mathbf C^{\mathbf O}$. Si par exemple $k$ est un anneau donné avec unité, on peut
considérer sur $X$ le faisceau d'anneaux constant correspondant, noté $k_X$.
La catégorie des $k_X$-Modules n'est autre que la catégorie des faisceaux de
$k$-modules sur $X$ ; si par exemple $k=\mathbf Z$, on obtient la catégorie
des faisceaux de groupes abéliens. Rappelons les faits suivants, signalés en
passant dans 1.5 et 1.9 :

\LayerTarget{tohoku:prop:3.1.1}
\textsc{Proposition 3.1.1.} \textit{Soit $\mathbf O$ un faisceau d'anneaux
avec unité sur l'espace $X$. Alors la catégorie additive $\mathbf C^{\mathbf O}$ des
$\mathbf O$-Modules sur $X$ est une catégorie abélienne satisfaisant les
axiomes \textup{AB 5)} et \textup{AB 3*)}, et admet un générateur.}

\LayerTarget{tohoku:proof:3.1.1:p1}
Signalons que la somme directe $\mathrm S$ d'une famille $(F_i)$ de faisceaux $F_i$ se
construit simplement en prenant pour chaque ouvert $U$ la somme directe des
$F_i(U)$ et passant au faisceau associé au préfaisceau ainsi construit ; le
procédé est le même pour la construction du produit $P$ des faisceaux $F_i$.
La différence essentielle dans les deux cas est que pour tout $x\in X$,
$\mathrm S(x)$ est bien la somme directe des $F_i(x)$, mais $P(x)$ \textit{n'est pas} le
produit direct des $F_i(x)$. On se convainc aussi aisément que l'axiome
AB 4*) n'est pas satisfait en général (en prenant par exemple pour
$\mathbf O$ le faisceau constant $\mathbf Z_X$). Enfin, rappelons que si pour
tout ouvert $U\subset X$, on désigne par $\mathbf O_U$ le faisceau de
$\mathbf O$-modules dont la restriction à $\complement U$ est nulle, et qui sur $U$
coïncide avec $\mathbf O$ [4, XVII, prop. 1], la famille des $\mathbf O_U$ est
un système de générateurs de $\mathbf C^{\mathbf O}$, comme il résulte facilement de
la définition 1.9 et du fait que les homomorphismes de $\mathbf O_U$

\TohokuContinuation{156}
